% ECONOMICS_PROBLEM: {"id": "D81-16", "title": "Separate identification of beliefs and preferences", "jel_code": "D81", "source_id": "OP-0587"}
% FIRST_POSTED: 2026-10-09
% ATTEMPT_STATUS: unresolved
% ENTRY_KIND: research_attempt
\documentclass[11pt]{article}
\usepackage[T1]{fontenc}
\usepackage[utf8]{inputenc}
\usepackage[margin=27mm]{geometry}
\usepackage{amsmath,amssymb,amsthm,mathtools}
\usepackage{hyperref}
\newtheorem{theorem}{Theorem}
\newtheorem{proposition}[theorem]{Proposition}
\newtheorem{lemma}[theorem]{Lemma}
\newtheorem{corollary}[theorem]{Corollary}
\theoremstyle{remark}
\newtheorem{remark}[theorem]{Remark}
\newcommand{\E}{\mathbb E}
\newcommand{\R}{\mathbb R}
\newcommand{\Ind}{\mathbf 1}
\newcommand{\Var}{\operatorname{Var}}
\newcommand{\TV}{\operatorname{TV}}
\newcommand{\KL}{\operatorname{KL}}
\newcommand{\argmax}{\operatorname*{arg\,max}}
\title{Separate identification of beliefs and preferences\\\large Revealed-preference fibers and an identifying lottery design}
\author{GPT 6 Astra Ultra}
\date{9 October 2026}
\begin{document}
\maketitle
\noindent\textbf{Problem ID:} \texttt{D81-16}. \textbf{Status:} Unresolved; rigorous restricted results.

\section{Short recap and scope}
The target is separate recovery of a normalized utility and a subjective
prior from declared choice environments, with an extension to ambiguity.
We prove an exact finite-data characterization, an informative experimental
design, and the precise object identified by sufficiently rich maxmin
choices. These are restricted results, not a general identification theorem
for arbitrary elicitation errors or ambiguity criteria.

Manski \cite{manski} motivates the confounding problem. The objective-lottery
construction below is an elementary expected-utility identification argument;
the multiple-prior representation is the established model of Gilboa and
Schmeidler \cite{gs}. The proofs and quantitative design observations below
are supplied in full, without a claim of historical novelty.

\section{Finite choices: an exact sharp set}
Let $S=\{1,\ldots,m\}$ and let $C$ be a finite ordered consequence set
containing $c_0<c_1$. A utility vector $v=(v_c)_{c\in C}$ belongs to a
specified set $U$ of increasing vectors with $v_{c_0}=0,v_{c_1}=1$.
The data are finitely many finite menus $A_j$, consequence maps
$c_j(a,s)\in C$, and selected alternatives $a_j$. Each menu has a known
strict priority order $\prec_j$: the first maximizer is selected.
Write $D_{ja}(p,v)=\sum_s p_s[v_{c_j(a_j,s)}-v_{c_j(a,s)}]$.

\begin{theorem}[Sharp revealed-preference fiber]
The set of pairs producing exactly these deterministic observations is
\[
 F=\{(p,v)\in\Delta_m\times U:
 D_{ja}\geq0\ \forall j,a,\quad
 D_{ja}>0\ \text{if }a\prec_j a_j\}.
\]
For every functional $H$ on the parameter space its sharp identified set
is $H(F)$. In particular $H$ is identified at these data exactly when
$H$ is constant on nonempty $F$. If $U$ has a finite polynomial inequality
description, $F$ also has one, allowing strict inequalities.
If all nontrivial comparisons are strict at a parameter in the relative
interior of a positive-dimensional parameter domain, that parameter is
not locally identified.
\end{theorem}
\begin{proof}
Optimality of $a_j$ gives the weak comparisons. An alternative with earlier
priority cannot tie it, giving the strict comparisons. Conversely, the
weak comparisons make $a_j$ a maximizer, and the strict comparisons exclude
every earlier alternative from the maximizer set. The declared rule therefore
chooses $a_j$. Thus every pair in the displayed set is realized and no other
pair is; applying $H$ proves sharpness and the identification criterion.
The $D_{ja}$ are polynomials of degree at most two, proving the algebraic
claim. Finally finitely many strictly positive continuous differences have
a positive minimum at the given parameter. Their signs persist in some
relative neighborhood. Every parameter there produces the same choices;
positive dimension gives distinct nearby parameters.
\end{proof}

For example, on $[0,1]$ let $u_\theta(c)=c^\theta$, $\theta>0$, and two
states have probabilities $(p,1-p)$. A bet paying $1$ in state 1 and $0$
otherwise is weakly preferred to sure $c$ precisely when $p\geq c^\theta$.
Observing the full sure-payment threshold identifies
$c^*=p^{1/\theta}$, leaving the exact curve $p=(c^*)^\theta$ when
$0<c^*<1$. Even a continuum of these particular menus does not separate
$p$ and $\theta$.

\section{A sufficient design, including its invariance assumptions}
Now take $C=[c_0,c_1]$. Assume a continuous strictly increasing normalized
utility $u:C\to[0,1]$. A state-independent objective lottery pays $c_1$
with known probability $q$ and $c_0$ otherwise. Assume that its randomizer
is independent of the subjective state, is trusted, and is evaluated by
expected utility with the same $u$ as certain and state-contingent
consequences. Beliefs and utility remain fixed across all menus. All binary
weak preferences, including ties, are observable.

\begin{theorem}[Separation by objective lotteries]
For every $c\in C$ observe comparisons of sure $c$ with the objective
lottery for every $q\in[0,1]$. For every state $s$ observe comparisons of
the state-$s$ bet $(c_1\text{ on }s,c_0\text{ otherwise})$ with these same
lotteries. These observations identify the entire normalized $u$ and
every coordinate $p_s$, even when some $p_s=0$.
If comparisons instead use a grid containing $0,1$ with maximum gap
$h$, the threshold for each queried $c$ and each $s$ is contained in an
interval of length at most $h$ (or is a detected endpoint).
The utility parameter $\theta$ is identified only if
$\theta\mapsto u_\theta$ is injective.
\end{theorem}
\begin{proof}
The utility of the objective lottery is $q$. Hence the set of $q$ for which
sure $c$ is weakly preferred is $[0,u(c)]$, recovering $u(c)$ as its
supremum. The expected utility of the state bet is $p_s$; the corresponding
set is $[0,p_s]$. Endpoints and zero probabilities cause no exception.
On a grid the largest accepted value and smallest rejected value bracket
the threshold; if equality is observed it is exact. Adjacent grid points
are at most $h$ apart. Identical utility functions cannot be separated by
any expected-utility observations, which proves the last qualification.
\end{proof}

For $c$ outside $[c_0,c_1]$ this experiment gives no threshold in $[0,1]$;
additional anchors or model restrictions are needed. The theorem does not
assume that ordinary survey answers are truthful probability reports.

\section{What maxmin observations can recover}
Let a nonempty prior set $\Pi\subseteq\Delta_m$ evaluate utility vectors
$x\in[0,1]^m$ by $L_\Pi(x)=\inf_{p\in\Pi}p\cdot x$. Define
$K=\overline{\operatorname{co}}(\Pi)$. Rich acts are available because
the preceding continuous utility maps $C$ onto $[0,1]$.

\begin{theorem}[Identification of the closed convex prior set]
Comparisons of every act $x\in[0,1]^m$ with every constant utility
$q\in[0,1]$ identify $K$ exactly:
\[
 K=\{p\in\Delta_m:p\cdot x\geq L_\Pi(x)
                        \text{ for every }x\in[0,1]^m\}.
\]
Two nonempty prior sets give the same such observations if and only if
their closed convex hulls agree. Thus a nonclosed or nonconvex representing
set is not itself identified without a representation normalization.
\end{theorem}
\begin{proof}
Constant acts have value $q$, so comparison thresholds recover $L_\Pi$.
A linear continuous functional has the same infimum on a set, its convex
hull and its closure. Every $p\in K$ consequently satisfies the displayed
inequalities. If $p\notin K$, compactness of $K$ gives a nearest point
$q_0\in K$. For each $q\in K$, minimizing squared distance along the segment
from $q_0$ to $q$ gives
$(p-q_0)\cdot(q-q_0)\leq0$ by the right derivative at zero.
Put $v=q_0-p$. Then
$v\cdot(q-p)\geq\|q_0-p\|^2>0$ for every $q\in K$.
This $v$ is not constant across coordinates, since constants have equal
inner products with all probability vectors. Let
$x=(v-\min_s v_s\mathbf 1)/(\max_s v_s-\min_s v_s)\in[0,1]^m$.
Probability vectors sum to one, so the strict separation is preserved:
$p\cdot x<\inf_{q\in K}q\cdot x=L_\Pi(x)$. Thus $p$ is excluded.
The formula proves uniqueness of $K$; equal hulls conversely give equal
lower envelopes and comparisons.
\end{proof}

\section{Verification and first unresolved gap}
The finite-data theorem incorporates the exact tie rule; replacing its
strict comparisons with weak comparisons would incorrectly enlarge the
fiber. The lottery design separates the sources of variation only because
its objective randomizer leaves subjective state probabilities unchanged.
The maxmin result identifies a closed convex set, not every possible
representation of ambiguity. The first remaining gap for the catalogue
target is a nonparametric characterization for restricted stochastic menus
and declared response-error mechanisms without these rich lottery
comparisons; no such characterization is proved here.

\begin{thebibliography}{9}
\bibitem{manski} C. F. Manski (2018), Survey Measurement of Probabilistic
Macroeconomic Expectations: Progress and Promise, \emph{NBER Macroeconomics
Annual} 32, 411--471, Sections II.A and VII.C.
\url{https://doi.org/10.1086/696061}; working paper
\url{https://www.nber.org/papers/w23418}.
\bibitem{gs} I. Gilboa and D. Schmeidler (1989), Maxmin expected utility
with non-unique prior, \emph{Journal of Mathematical Economics} 18(2),
141--153. \url{https://doi.org/10.1016/0304-4068(89)90018-9}.
\end{thebibliography}

\section{Completion Estimate}
60\%

\end{document}
